\documentclass[11pt]{article}
\usepackage[margin=1in]{geometry}
\usepackage{amsmath,amssymb,amsthm}
\usepackage[colorlinks=true,urlcolor=blue,linkcolor=blue]{hyperref}
\newcommand{\D}{\mathbb D}
\newcommand{\C}{\mathbb C}
\newtheorem{theorem}{Theorem}
\newtheorem{proposition}[theorem]{Proposition}
\title{A cubic symmetric curve avoiding the diagonal\\
\large A solution to MathDB 398865}
\author{Principia Math in Progress}
\date{September 12, 2026}
\begin{document}
\maketitle

\paragraph{Task 1. Construct the requested polynomial.}
Let $\D=\{z\in\C:|z|<1\}$. The problem asks for a symmetric polynomial
$P(z_1,z_2)$ whose zero set in $\D^2$ is nonempty and connected, but does
not meet $\{(z,z):z\in\D\}$. This is the two-variable construction requested
by Khavinson in \cite[p.~381]{source}. The answer has the least
possible total degree.

\begin{theorem}
The polynomial
\[
\boxed{\quad P(z_1,z_2)=(z_1+z_2)(z_1-z_2)^2-\frac12 \quad}
\]
is symmetric. Its zero set $V=\{P=0\}\cap\D^2$ is nonempty and
path connected, and it does not meet the diagonal, even in $\C^2$.
\end{theorem}

\begin{proof}
Symmetry follows from the square, and $P(z,z)=-1/2$.
Introduce coordinates
\[
u=\frac{z_1+z_2}{2},\qquad w=\frac{z_1-z_2}{2}.
\]
The equation $P=0$ becomes $uw^2=1/16$. In particular $w\ne0$.
Set $a=1/16$. The map
\[
\Phi(w)=\left(\frac{a}{w^2}+w,\frac{a}{w^2}-w\right),\qquad w\in\C\setminus\{0\},
\]
is a bijective parametrization of the entire affine zero set, with
continuous inverse $w=(z_1-z_2)/2$. Consequently $V$ is homeomorphic to
\[
\Omega=\left\{w\ne0:
 \left|\frac{a}{w^2}+w\right|<1,\quad
 \left|\frac{a}{w^2}-w\right|<1\right\}.
\]
It remains to establish connectedness of this entire parameter set.

Write $w=re^{i\theta}$ and $c=|\cos(3\theta)|\in[0,1]$. Direct expansion gives
\begin{equation}\label{eq:radial}
\max\left\{\left|\frac{a}{w^2}+w\right|^2,
           \left|\frac{a}{w^2}-w\right|^2\right\}
=F_c(r):=r^2+\frac{a^2}{r^4}+\frac{2ac}{r}.
\end{equation}
For every fixed $c\in[0,1]$ and $r>0$,
\[
F_c''(r)=2+\frac{20a^2}{r^6}+\frac{4ac}{r^3}>0.
\]
Thus $\{r>0:F_c(r)<1\}$ is an interval. Moreover,
\[
F_c(1/2)=\frac{5+4c}{16}\le\frac9{16}<1.
\]
The circle $|w|=1/2$ therefore lies wholly in $\Omega$. If
$re^{i\theta}\in\Omega$, the radial segment
\[
s\longmapsto\bigl((1-s)r+s/2\bigr)e^{i\theta},
\qquad 0\le s\le1,
\]
also lies in $\Omega$, by convexity of $F_c$. Any two points of $\Omega$
can now be joined by following their radial segments to that circle and
then following the circle. This proves path connectedness and nonemptiness.
An explicit zero is $\Phi(1/2)=(3/4,-1/4)$.
\end{proof}

The radial formula is a strong deformation retraction of $\Omega$
onto $|w|=1/2$: written as
$R_s(w)=((1-s)+s/(2|w|))w$, it is continuous without choosing an argument.
It transports through $\Phi$ to a deformation retraction of $V$ onto a
circle. In particular the construction does not assert simple connectedness.
The affine curve is smooth, since the derivative with respect to $u$ of
$8uw^2-1/2$ is nonzero on it.

\begin{proposition}
No symmetric polynomial of total degree at most two has a nonempty,
connected zero set in $\D^2$ that avoids the diagonal.
\end{proposition}

\begin{proof}
In the same coordinates, any such polynomial has the form
$A w^2+Q(u)$, where $A\in\C$ and $\deg Q\le2$.
Diagonal avoidance implies $Q(u)\ne0$ for every $u\in\D$.
If $A=0$, there are no zeros in $\D^2$, because
$|u|\le(|z_1|+|z_2|)/2<1$.

If $A\ne0$, the nonvanishing holomorphic function $-Q/A$ has a
nonvanishing holomorphic square root $g$ on the simply connected disk.
Every zero in $\D^2$ belongs to one of the disjoint sets $w=g(u)$ and
$w=-g(u)$. The continuous function $w/g(u)$ on the zero set takes values
in $\{1,-1\}$, so these sets are both relatively open and closed.
If the zero set is nonempty, both signs occur: interchanging $z_1,z_2$
preserves the bidisk and replaces $w$ by $-w$. Thus the zero set is
disconnected.
\end{proof}

\paragraph{Scope and verification.}
The argument treats all zeros in the open bidisk, rather than a selected
component or a finite sample. Separate linearity belongs to the preceding
Walsh theorem in the source, and is not imposed on the two-variable request.
The source's three-variable example is background and receives no credit here.
Symbolic identities and exact rational path checks are saved alongside
this note; the proofs above establish the universal claims.

\paragraph{Task 2. Classify cubic examples.}
Characterize the symmetric polynomials of total degree three whose
nonempty zero set in the bidisk is connected and avoids the diagonal.
Identify useful normal forms and the role of zeros of the coefficient
of $w^2$ in the $(u,w)$ coordinates.

\paragraph{Task 3. Determine the topology of the parameter family.}
For $P_c(z_1,z_2)=(z_1+z_2)(z_1-z_2)^2-c$, determine every value
$c\in\C\setminus\{0\}$ for which the zero set in $\D^2$ is nonempty and
connected. Determine the number and topology of components at the
remaining parameter values and at transition values.

\paragraph{Task 4. Quantify the conformal geometry.}
For the connected curves in Task 3, determine their annulus moduli,
or sharp asymptotics for those moduli as $c$ tends to zero and as the
connected region approaches a topology transition.

\paragraph{Task 5. Avoid all coordinate coincidences in more variables.}
For each $N\ge3$, construct, or rule out, a symmetric polynomial whose
nonempty zero set in $\D^N$ is connected and avoids every hyperplane
$z_i=z_j$ with $i\ne j$. Determine the least possible total degree when
such a polynomial exists.

\begin{thebibliography}{9}
\bibitem{source}
D.~Khavinson, \emph{Problems for Polynomials of One and several Variables},
in \emph{Dirichlet Series and Function Theory in Polydiscs},
Oberwolfach Reports 2014/6, pp.~381--383, specifically p.~381.
\url{https://ems.press/journals/owr/articles/12866}.
\bibitem{mathdb}
MathDB 398865, \emph{Symmetric polynomial zero sets in the bidisk}.
\url{https://mathdb.com/p/398865}. Live status and public-claim checks
archived September 12, 2026.
\end{thebibliography}
\end{document}
